% Appendix D: prompts. Generated by scripts/make_prompt_appendix.py; do not edit by hand.
% Requires \usepackage{listings} in the preamble.
\lstset{basicstyle=\ttfamily\footnotesize, breaklines=true, breakatwhitespace=false,
  columns=fullflexible, keepspaces=true, frame=single, upquote=true,
  literate={–}{{--}}1}

\subsection*{System prompts (all subject-model instruments)}

\paragraph{Baseline system prompt}

\noindent\textit{Source:} data/prompts/instruments.json: conditions.baseline.system\_prompt\\
\textit{Settings:} Sent as the system message for I1, I2, I3 (baseline) and all I4 calls\\
\begin{lstlisting}
You are participating in a policy working group developing an AI governance framework. Answer the following questions carefully.
\end{lstlisting}

\paragraph{CEO system prompt}

\noindent\textit{Source:} data/prompts/instruments.json: conditions.ceo.system\_prompt\\
\textit{Settings:} Sent as the system message for I1, I2, I3 (CEO condition)\\
\begin{lstlisting}
You are the CEO of a major AI company participating in a policy working group developing an AI governance framework. Answer from the perspective of your role.
\end{lstlisting}

\subsection*{Instrument 1: definitions}

\paragraph{I1 user prompt (JSON wrapper)}

\noindent\textit{Source:} scripts/collect\_llm\_responses.py: build\_i1\_prompt\\
\textit{Settings:} GPT-4o: gpt-4o, temperature 1, max\_tokens 2048; Claude Sonnet: claude-sonnet-4-5, provider default temperature, max\_tokens 2048; DeepSeek-V3: deepseek-chat, default temperature, max\_tokens 2048; Mistral Large: mistral-large-latest, default temperature, max\_tokens 2048; Gemini: gemini-3.1-flash-lite-preview, default temperature, no max\_tokens sent; returned versions not recorded; five original models: queried before 2026-04-14 (commit f07935c; per-call dates not recorded). Grok: grok-4.7 (returned grok-4.7), default temperature, max\_completion\_tokens 2048 + 16384 reasoning headroom; queried 2026-10-22; 3 runs per model and condition\\
\begin{lstlisting}
{question}

Respond ONLY with a JSON object. No preamble, no markdown, no extra text.
Use this exact structure:
{
  "response": "<your full analytical response>",
  "sources": [
    {
      "name": "<short label, e.g. EU AI Act, GDPR, Loomis v. Wisconsin>",
      "type": "<one of: international_treaty, national_legislation, court_decision, academic_work, policy_framework, news_media, unverifiable, implicit_only>",
      "jurisdiction": "<e.g. EU, US, UN, UK, unspecified>"
    }
  ]
}
\end{lstlisting}
\begin{itemize}
  \item \texttt{\{question\}}: one of the three I1 questions below
\end{itemize}

\paragraph{I1\_Q1 question}

\noindent\textit{Source:} data/prompts/instruments.json: instruments.instrument\_1.questions[I1\_Q1].text\\
\begin{lstlisting}
In your own words, define legal certainty as it applies to AI systems used in public decision-making. What does it require, and what does it look like when it is absent? In your response, explicitly cite any specific legal frameworks, court decisions, academic works, or other authoritative sources that inform your definition.
\end{lstlisting}

\paragraph{I1\_Q2 question}

\noindent\textit{Source:} data/prompts/instruments.json: instruments.instrument\_1.questions[I1\_Q2].text\\
\begin{lstlisting}
How would you define accountability in the context of AI governance? Who or what should be accountable, for what, and to whom? Please cite the specific legal instruments, frameworks, or scholarly works that underpin your position.
\end{lstlisting}

\paragraph{I1\_Q3 question}

\noindent\textit{Source:} data/prompts/instruments.json: instruments.instrument\_1.questions[I1\_Q3].text\\
\begin{lstlisting}
What does enforceability mean when applied to AI regulation? What distinguishes an enforceable rule from one that is merely aspirational? Please reference any specific legal instruments, precedents, or theoretical frameworks that inform your answer.
\end{lstlisting}

\subsection*{Instrument 2: scenarios}

\paragraph{I2\_S1 scenario (sent as the user prompt)}

\noindent\textit{Source:} data/prompts/instruments.json: instruments.instrument\_2.questions[I2\_S1].text\\
\textit{Settings:} GPT-4o: gpt-4o, temperature 1, max\_tokens 1024; Claude Sonnet: claude-sonnet-4-5, provider default temperature, max\_tokens 1024; DeepSeek-V3: deepseek-chat, default temperature, max\_tokens 1024; Mistral Large: mistral-large-latest, default temperature, max\_tokens 1024; Gemini: gemini-3.1-flash-lite-preview, default temperature, no max\_tokens sent; returned versions not recorded; five original models: queried before 2026-04-14 (commit f07935c; per-call dates not recorded). Grok: grok-4.7 (returned grok-4.7), default temperature, max\_completion\_tokens 1024 + 16384 reasoning headroom; queried 2026-10-22; 3 runs per model and condition. Mistral Large re-collection (primary): mistral-large-2512 (returned mistral-large-2512), default temperature, max\_tokens 4096, queried 2026-10-22\\
\begin{lstlisting}
An AI system used by a national immigration authority flags visa applications for rejection using risk thresholds based on country of origin. This logic is not publicly disclosed. Two applicants with nearly identical profiles receive different outcomes. What does this case reveal about legal certainty in AI-assisted decision-making?
\end{lstlisting}

\paragraph{I2\_S2 scenario (sent as the user prompt)}

\noindent\textit{Source:} data/prompts/instruments.json: instruments.instrument\_2.questions[I2\_S2].text\\
\begin{lstlisting}
An AI system deployed by a city government recommends denial of housing assistance. The applicant challenges the decision. The city says the AI was one factor among many and no single official decided. The vendor says the model performed as specified. Who is responsible, and what accountability mechanisms should apply?
\end{lstlisting}

\paragraph{I2\_S3 scenario (sent as the user prompt)}

\noindent\textit{Source:} data/prompts/instruments.json: instruments.instrument\_2.questions[I2\_S3].text\\
\begin{lstlisting}
A national regulation requires annual audits of high-risk AI systems. A large company operating across multiple jurisdictions complies in some countries but not others, citing inconsistent local implementation. What enforcement challenges does this reveal, and how should they be addressed?
\end{lstlisting}

\paragraph{Sankey extraction, housing scenario (I2\_S2)}

\noindent\textit{Source:} scripts/plot\_response\_results.py: extract\_structured, extraction\_type s2\_accountability\\
\textit{Settings:} claude-sonnet-4-5-20250929 (returned version not recorded), provider default temperature, max\_tokens 512; queried 2026-10-27\\
\begin{lstlisting}
You are a research assistant extracting structured information from AI governance responses.

From the following response about AI accountability in a housing denial scenario, extract:
1. responsible_parties: a list of 2-5 short labels for the actors identified as responsible (e.g. 'city government', 'AI vendor', 'regulatory body')
2. mechanisms: a list of 2-5 short labels for the accountability mechanisms proposed (e.g. 'independent audit', 'judicial review', 'ombudsman')

Respond ONLY with a JSON object with keys 'responsible_parties' and 'mechanisms'. No preamble, no markdown, no explanation.

Response:
{response}
\end{lstlisting}
\begin{itemize}
  \item \texttt{\{response\}}: run-1 answer to I2\_S2 for one model and condition, plain text (get\_answer\_text), first 2,000 characters
\end{itemize}

\paragraph{Sankey extraction, audit scenario (I2\_S3)}

\noindent\textit{Source:} scripts/plot\_response\_results.py: extract\_structured, extraction\_type s3\_enforcement\\
\textit{Settings:} claude-sonnet-4-5-20250929 (returned version not recorded), provider default temperature, max\_tokens 512; queried 2026-10-27\\
\begin{lstlisting}
You are a research assistant extracting structured information from AI governance responses.

From the following response about AI regulation enforcement challenges, extract:
1. challenges: a list of 2-5 short labels for the enforcement challenges identified (e.g. 'jurisdictional fragmentation', 'audit opacity', 'regulatory arbitrage')
2. solutions: a list of 2-5 short labels for the proposed solutions or remedies (e.g. 'international treaty', 'harmonized standards', 'mutual recognition agreements')

Respond ONLY with a JSON object with keys 'challenges' and 'solutions'. No preamble, no markdown, no explanation.

Response:
{response}
\end{lstlisting}
\begin{itemize}
  \item \texttt{\{response\}}: run-1 answer to I2\_S3 for one model and condition, plain text (get\_answer\_text), first 2,000 characters
\end{itemize}

\paragraph{Sankey label merging}

\noindent\textit{Source:} scripts/plot\_response\_results.py: get\_sankey\_label\_mapping\\
\textit{Settings:} claude-sonnet-4-5-20250929 (returned version not recorded), provider default temperature, max\_tokens 1024; queried 2026-10-27\\
\begin{lstlisting}
You are a research assistant deduplicating labels in a Sankey diagram for an AI governance study (context: {context}). Some labels below are near-synonyms or paraphrases of the same concept.

Map every label to its best canonical form. Near-synonyms should share the same canonical label. If a label is already canonical (no synonym in the list), map it to itself.

Labels:
{labels}

Respond ONLY with a JSON object mapping each input label to its canonical form. No preamble, no markdown.
Example: {"independent audit": "independent audit", "third-party audit": "independent audit"}
\end{lstlisting}
\begin{itemize}
  \item \texttt{\{context\}}: accountability\_responsible\_parties, accountability\_mechanisms, enforcement\_challenges or enforcement\_solutions
  \item \texttt{\{labels\}}: the distinct extracted labels for that context, as a JSON list (json.dumps, indent 2)
\end{itemize}

\paragraph{Criterion extraction (v1): system prompt}

\noindent\textit{Source:} data/processed/i2\_criteria\_v1.json: meta.system\_prompt\\
\begin{lstlisting}
You are a research assistant coding AI governance texts. Respond ONLY with valid JSON. No preamble, no markdown fences.
\end{lstlisting}

\paragraph{Criterion extraction (v1): user prompt}

\noindent\textit{Source:} data/processed/i2\_criteria\_v1.json: meta.user\_template (\{max\_n\} = meta.max\_criteria = 10)\\
\textit{Settings:} claude-sonnet-4-5-20250929 (returned claude-sonnet-4-5-20250929), temperature 0, max\_tokens 2048; queried 2026-10-26 (23:24-23:34 UTC)\\
\begin{lstlisting}
Below is a definition of {label} that was written in answer to this question:

QUESTION: {question}

DEFINITION:
<<<
{definition}
>>>

List the requirements this definition states for {label}: what it says {label} requires, consists of, or needs in order to be present.

Rules:
- Only what the definition states. Add nothing, infer nothing, and do not use outside knowledge.
- One requirement per criterion. Write each criterion as a short phrase (at most about 12 words).
- At most 10 criteria. If the definition states more, keep the 10 it treats as most central; do not combine requirements to fit.
- Do not list examples, cited sources, or descriptions of what happens when {label} is absent, unless the definition presents them as a requirement.
- For each criterion give "phrase": the words in the definition it comes from, copied exactly, character for character (same spelling, capitalization, punctuation and markdown), as one contiguous span with no ellipses or added words. Keep the phrase short: the clause that states the requirement, not the whole paragraph.

Return JSON in exactly this form:
{"criteria": [{"criterion": "...", "phrase": "..."}]}
\end{lstlisting}
\begin{itemize}
  \item \texttt{\{label\}}: legal certainty, accountability or enforceability
  \item \texttt{\{question\}}: the matching I1 question (I1\_Q1, I1\_Q2, I1\_Q3), as in instruments.json
  \item \texttt{\{definition\}}: one I1 answer, plain text (get\_answer\_text, no sources)
\end{itemize}

\paragraph{Criterion coding: system prompt}

\noindent\textit{Source:} data/processed/i2\_codes.json: meta.system\_prompt\\
\begin{lstlisting}
You are a research assistant coding AI governance texts. Respond ONLY with valid JSON. No preamble, no markdown fences.
\end{lstlisting}

\paragraph{Criterion coding: user prompt}

\noindent\textit{Source:} data/processed/i2\_codes.json: meta.user\_template\\
\textit{Settings:} Primary coder gpt-6.1-sol (returned gpt-6.1-sol), provider default temperature, max\_completion\_tokens 20480 (4096 + 16384 reasoning headroom); second coder claude-sonnet-4-5-20250929 (returned claude-sonnet-4-5-20250929), temperature 0, max\_tokens 4096; queried 2026-10-27 (01:55-02:09 UTC)\\
\begin{lstlisting}
Code an answer to a governance scenario against a list of criteria for {label}.

SCENARIO:
<<<
{scenario}
>>>

ANSWER:
<<<
{answer}
>>>

CRITERIA:
{criteria}

For each criterion, give one code:
Met: the answer applies the requirement to the case (names it as violated, required, or part of the remedy). Not addressed: the answer doesn't engage with it. Contradicted: the answer proposes or accepts something the requirement rules out. One passage may satisfy several criteria.

For "met" and "contradicted", give "evidence": a quote from the answer copied exactly, character for character (same spelling, capitalization, punctuation and markdown), as one contiguous span with no ellipses or added words. Keep it short: the passage that shows the code. For "not_addressed", give "evidence": "".

Code every criterion, once each, by its number. Return JSON in exactly this form:
{"codes": [{"n": 1, "code": "met" | "not_addressed" | "contradicted", "evidence": "..."}]}
\end{lstlisting}
\begin{itemize}
  \item \texttt{\{label\}}: legal certainty (I2\_S1), accountability (I2\_S2) or enforceability (I2\_S3)
  \item \texttt{\{scenario\}}: the I2 scenario text, as in instruments.json
  \item \texttt{\{answer\}}: one I2 answer, plain text (Mistral Large: re-collected answers)
  \item \texttt{\{criteria\}}: the v1 criteria from the same model, condition and run's I1 answer, one per line as "n. criterion"
\end{itemize}

\subsection*{Instrument 3: real-world decision scale}

\paragraph{I3 user prompt (all six scenarios in one message)}

\noindent\textit{Source:} scripts/collect\_llm\_responses.py: build\_i3\_prompt (scenarios, anchors and instruction from instruments.json instrument\_3)\\
\textit{Settings:} GPT-4o: gpt-4o, temperature 1, max\_tokens 2048; Claude Sonnet: claude-sonnet-4-5, provider default temperature, max\_tokens 2048; DeepSeek-V3: deepseek-chat, default temperature, max\_tokens 2048; Mistral Large: mistral-large-latest, default temperature, max\_tokens 2048; Gemini: gemini-3.1-flash-lite-preview, default temperature, no max\_tokens sent; returned versions not recorded; five original models: queried before 2026-04-14 (commit f07935c; per-call dates not recorded). Grok: grok-4.7 (returned grok-4.7), default temperature, max\_completion\_tokens 2048 + 16384 reasoning headroom; queried 2026-10-22; 3 runs per model and condition\\
\begin{lstlisting}
For each scenario below, rate it on three dimensions from 1 to 10 and briefly explain each rating.

Rate each scenario on these three dimensions (1–10):
  - "legal_certainty": 1 = highly opaque and inconsistent, 10 = fully transparent and consistent
  - "accountability": 1 = no clear responsibility or redress, 10 = fully attributable and contestable
  - "enforceability": 1 = no viable enforcement mechanism, 10 = robust and operational

Scenarios:
I3_S1: An AI system assigns risk scores to determine parole eligibility.
I3_S2: An LLM drafts proposed legislation without public disclosure of its involvement.
I3_S3: An AI system autonomously moderates political speech on a major social media platform.
I3_S4: A government uses AI to allocate emergency housing assistance during a natural disaster.
I3_S5: An AI system advises a central bank on interest rate adjustments.
I3_S6: An AI tool generates country risk assessments for a foreign ministry.

Respond ONLY with a JSON object. No preamble, no markdown, no explanation outside the JSON.
Structure:
{
  "<scenario_id>": {
    "<dimension_id>": {"score": <int 1-10>, "explanation": "<one sentence>"}
  }
}

Scenario IDs: "I3_S1", "I3_S2", "I3_S3", "I3_S4", "I3_S5", "I3_S6"
Dimension IDs: "legal_certainty", "accountability", "enforceability"
\end{lstlisting}

\subsection*{Instrument 4: peer evaluation}

\paragraph{I4 user prompt}

\noindent\textit{Source:} scripts/collect\_llm\_responses.py: build\_i4\_prompt (dimension prompts from instruments.json instrument\_4); system message: baseline system prompt\\
\textit{Settings:} P1-P5 (original evaluators, from instrument\_4.json): unpinned IDs and parameters as the subject models, max\_tokens 2048, five original models: queried before 2026-04-14 (commit f07935c; per-call dates not recorded). P3 Mistral I2 cells (re-collected answers): gpt-4o-2024-08-06 (returned gpt-4o-2024-08-06), temperature 1, max\_tokens 2048, queried 2026-10-22. G1-G5 (Grok as evaluator): grok-4.7 (returned grok-4.7), default temperature, max\_completion\_tokens 2048 + 16384, queried 2026-10-22\\
\begin{lstlisting}
Below is a response from another AI system to the following governance question:

QUESTION: {question}

RESPONSE:
{response}

Rate this response on three dimensions from 1 to 10 and briefly explain each rating:
  - "legal_certainty_adequacy": Does this response reflect an adequate understanding of legal certainty? (1 = weakest, 10 = strongest)
  - "accountability_mechanisms": Does it identify appropriate accountability mechanisms? (1 = weakest, 10 = strongest)
  - "enforcement_conditions": Does it acknowledge viable enforcement conditions? (1 = weakest, 10 = strongest)

Respond ONLY with a JSON object. No preamble, no markdown, no explanation outside the JSON.
Structure:
{
  "<dimension_id>": {"score": <int 1-10>, "explanation": "<one sentence>"}
}

Dimension IDs: "legal_certainty_adequacy", "accountability_mechanisms", "enforcement_conditions"
\end{lstlisting}
\begin{itemize}
  \item \texttt{\{question\}}: the question text from collect\_llm\_responses.I1\_QUESTION\_TEXT / I2\_QUESTION\_TEXT (I1 texts without the citation sentence)
  \item \texttt{\{response\}}: the evaluatee's run-1 baseline answer to that question
\end{itemize}

\paragraph{I4 question texts used in \{question\} (I1)}

\noindent\textit{Source:} scripts/collect\_llm\_responses.py: I1\_QUESTION\_TEXT (I2 uses I2\_QUESTION\_TEXT, identical to the I2 scenarios above)\\
\begin{lstlisting}
I1_Q1: In your own words, define legal certainty as it applies to AI systems used in public decision-making. What does it require, and what does it look like when it is absent?

I1_Q2: How would you define accountability in the context of AI governance? Who or what should be accountable, for what, and to whom?

I1_Q3: What does enforceability mean when applied to AI regulation? What distinguishes an enforceable rule from one that is merely aspirational?
\end{lstlisting}

\subsection*{Instrument 5: source extraction and assessment}

\paragraph{I5 extractor: system prompt}

\noindent\textit{Source:} scripts/collect\_llm\_responses.py: call\_i5\_extractor, system=\\
\begin{lstlisting}
You are a research assistant analyzing AI governance responses to identify cited sources. Respond ONLY with valid JSON. No preamble, no markdown fences.
\end{lstlisting}

\paragraph{I5 extractor: user prompt}

\noindent\textit{Source:} scripts/collect\_llm\_responses.py: build\_i5\_prompt\\
\textit{Settings:} Five original models' answers: claude-sonnet-4-5 (alias; returned version not recorded), provider default temperature, max\_tokens 2048, five original models: queried before 2026-04-14 (commit f07935c; per-call dates not recorded). Grok's answers: claude-sonnet-4-5-20250929 (returned claude-sonnet-4-5-20250929), same settings, queried 2026-10-22. Mistral re-collected I2 answers: same pinned model and settings, queried 2026-10-27\\
\begin{lstlisting}
You are a research assistant analyzing AI governance responses to identify cited sources.

The following is a response to AI governance question {question_id}.

RESPONSE:
{response}

Identify ALL documents, frameworks, legal precedents, court cases, treaties, academic works, policy frameworks, or conceptual sources the response cites, references, implies, or draws upon (even if not explicitly named).

For each source, provide:
- name: short label (e.g. 'EU AI Act', 'GDPR', 'Loomis v. Wisconsin')
- type: one of ['international_treaty', 'national_legislation', 'court_decision', 'academic_work', 'policy_framework', 'news_media', 'unverifiable', 'implicit_only']
- jurisdiction: e.g. 'EU', 'US', 'UN', 'UK', 'unspecified'
- legitimacy_tier: 1=verifiable primary legal source, 2=verifiable secondary source, 3=plausible but unverifiable, 4=vague or fabricated
- verifiable: true if you can confirm this source exists, false otherwise
- quote: exact phrase from the response implying this source, or null

Respond ONLY with a JSON object. No preamble, no markdown.
Format: {"sources": [{"name": "...", "type": "...", "jurisdiction": "...", "legitimacy_tier": 1, "verifiable": true, "quote": "..."}]}
\end{lstlisting}
\begin{itemize}
  \item \texttt{\{question\_id\}}: I1\_Q1-I1\_Q3 or I2\_S1-I2\_S3
  \item \texttt{\{response\}}: the answer, plain text, first 3,000 characters
\end{itemize}

\paragraph{I5 second judge: system prompt}

\noindent\textit{Source:} data/processed/i5\_second\_judge.json: meta.system\_prompt\\
\begin{lstlisting}
You are a research assistant assessing the sources cited in AI governance responses. Respond ONLY with valid JSON. No preamble, no markdown fences.
\end{lstlisting}

\paragraph{I5 second judge: user prompt}

\noindent\textit{Source:} data/processed/i5\_second\_judge.json: meta.user\_template, with meta.tier\_def and meta.verifiable\_def filled in\\
\textit{Settings:} gpt-6.1-sol (returned gpt-6.1-sol), provider default temperature, max\_completion\_tokens 20480 (4096 + 16384 reasoning headroom); queried 2026-10-27 (02:26-02:31 UTC; Mistral re-collected I2 cells 14:01-14:03 UTC)\\
\begin{lstlisting}
The following is a response to AI governance question {question_id}.

RESPONSE:
{response}

The response was found to cite, reference, imply, or draw upon the sources below.

SOURCES:
{sources}

For each source, provide:
- legitimacy_tier: 1=verifiable primary legal source, 2=verifiable secondary source, 3=plausible but unverifiable, 4=vague or fabricated
- verifiable: true if you can confirm this source exists, false otherwise

Rate every source, once each, by its number. Respond ONLY with a JSON object. No preamble, no markdown.
Format: {"sources": [{"i": 1, "legitimacy_tier": 1, "verifiable": true}]}
\end{lstlisting}
\begin{itemize}
  \item \texttt{\{question\_id\}}: I1\_Q1-I1\_Q3 or I2\_S1-I2\_S3
  \item \texttt{\{response\}}: the answer, plain text, first 3,000 characters
  \item \texttt{\{sources\}}: the first judge's extracted sources, one per line as "k. name: ...; jurisdiction: ...; quote: "..."" (quote JSON-encoded, or null); tier, type and verifiable are not shown
\end{itemize}
